% =============================================================================
%  Memoria del TFG — Frontend del Sistema de Corrección Digital de Exámenes Escritos
%  Autor: Alejandro Soler Sichling
%  Tutor: PENDIENTE (ver preamble/metadata.tex)
%  Grado en Ingeniería Informática — EPS, Universidad Autónoma de Madrid
%  Curso: 2026/2027
% =============================================================================
%  ESTADO: ESQUELETO. Estructura + material factual verificado + marcas \pendiente.
%  No se ha compilado: no había distribución LaTeX en la máquina de desarrollo.
%  Ver README.md (origen de la plantilla y receta de compilación).
%
%  Clase: tfgtfmthesisuam (E. Anguiano Rey, EPS-UAM), la misma que usa la memoria
%  del TFG de backend (repositorio SCDEE, docs/main.tex), para que ambas memorias
%  sean coherentes. Según la Normativa de TFG de la EPS (8/5/2026, art. 5.1.1) la
%  memoria sigue "las directrices establecidas en la plataforma de TFG": CONFIRMAR
%  en tfg.eps.uam.es / Moodle que esta plantilla es la exigida para el curso 26/27.
%
%  Receta (igual que la memoria del backend):
%    pdflatex main ; bibtex main ; makeglossaries main ;
%    makeindex -s main.ist main ; pdflatex main ; pdflatex main
% =============================================================================

\documentclass[epsbased,tfg,covers,final,printable,extendedindex,%
               firstnumbered,lof,lot,loc]{tfgtfmthesisuam}

% =============================================================================
%  Metadatos del documento (mismas variables que la memoria del backend)
% =============================================================================
\title{Frontend del Sistema de Corrección Digital de Exámenes Escritos}
% \subtitle{...}  % \pendiente: decidir si el subtítulo menciona la generación agéntica
\author{Alejandro Soler Sichling}
\advisor{PENDIENTE: nombre del tutor académico}
\levelin{Grado en Ingeniería Informática}
\copyrightdate{2026}

\coverdata{%
  Escuela Politécnica Superior\\
  Universidad Autónoma de Madrid\\
  C\textbackslash\ Francisco Tomás y Valiente nº 11%
}

\resumenfile{front-matter/resumen}
\abstractfile{front-matter/abstract}
% \ackfile{front-matter/agradecimientos}
\keywords{exam grading, frontend, React, API contract, coding agents, accessibility}
\palabrasclave{corrección de exámenes, frontend, React, contrato de API, agentes de programación, accesibilidad}

\graphicsdir{files}
\logosdir{logos}
\codesdir{codes}
\datadir{data}

\bibliographyconfig[ieeetr]{bibliography/references}

\pdfmetavalues
  {Alejandro Soler Sichling}
  {Frontend del Sistema de Corrección Digital de Exámenes Escritos}%
  {Trabajo Fin de Grado}%
  {corrección digital, exámenes, frontend, React, agentes, accesibilidad}

% =============================================================================
%  Acrónimos (convención de la memoria del backend: \ac, \acs, \acl)
% =============================================================================
\newacronym{api}{API}{Application Programming Interface}
\newacronym{css}{CSS}{Cascading Style Sheets}
\newacronym{eps}{EPS}{Escuela Politécnica Superior}
\newacronym{http}{HTTP}{Hypertext Transfer Protocol}
\newacronym{ia}{IA}{Inteligencia Artificial}
\newacronym{json}{JSON}{JavaScript Object Notation}
\newacronym{jwt}{JWT}{JSON Web Token}
\newacronym{llm}{LLM}{Large Language Model}
\newacronym{lms}{LMS}{Learning Management System}
\newacronym{mvp}{MVP}{Minimum Viable Product}
\newacronym{pdf}{PDF}{Portable Document Format}
\newacronym{rest}{REST}{Representational State Transfer}
\newacronym{spa}{SPA}{Single Page Application}
\newacronym{tfg}{TFG}{Trabajo Fin de Grado}
\newacronym{uam}{UAM}{Universidad Autónoma de Madrid}
\newacronym{ui}{UI}{User Interface}
\newacronym{ux}{UX}{User Experience}
\newacronym{wcag}{WCAG}{Web Content Accessibility Guidelines}

% =============================================================================
%  Definiciones (\newdefinition de la clase) y macros de redacción
% =============================================================================
\usepackage{xcolor}
% Marca visible para lo que falta redactar o decidir. Eliminar todas antes de entregar:
%   grep -rn "pendiente" chapters appendices front-matter
\newcommand{\pendiente}[1]{\textcolor{red}{\textbf{[PENDIENTE:} #1\textbf{]}}}

\newdefinition{backend}{backend}{backends}%
  {Parte del sistema responsable de la lógica de negocio, la persistencia y la
   exposición de servicios, sin contacto directo con la interfaz de usuario}
\newdefinition{frontend}{frontend}{frontends}%
  {Parte del sistema responsable de la interacción con el usuario y de la
   presentación de la información}
\newdefinition{instancia}{instancia de examen}{instancias de examen}%
  {Examen escaneado de un estudiante concreto, con sus páginas, anotaciones y notas}
\newdefinition{rubrica}{rúbrica}{rúbricas}%
  {Conjunto de criterios con puntuación con los que se califica un problema}
\newdefinition{mock}{simulación del backend}{simulaciones del backend}%
  {Implementación en el navegador que responde a las peticiones con los mismos
   formatos que el backend real, usada para desarrollar y probar sin él}
\newdefinition{agente}{agente de programación}{agentes de programación}%
  {Sistema basado en un modelo de lenguaje que lee y modifica código, ejecuta
   órdenes y verifica resultados con un grado de autonomía configurable}
\newdefinition{worktree}{worktree}{worktrees}%
  {Directorio de trabajo adicional de Git asociado a otra rama del mismo
   repositorio, que permite trabajar en paralelo sin compartir ficheros}

% Fragmentos de código (JavaScript/JSX) en el capítulo de implementación
\usepackage{listings}
\lstdefinelanguage{JavaScript}{%
  keywords={const,let,var,function,return,if,else,for,of,new,async,await,try,catch,
            finally,throw,export,import,from,default,null,true,false,typeof},
  sensitive=true,
  comment=[l]{//},
  morecomment=[s]{/*}{*/},
  morestring=[b]',
  morestring=[b]",
  morestring=[b]`}
\lstset{language=JavaScript, basicstyle=\ttfamily\footnotesize,
        keywordstyle=\bfseries, commentstyle=\itshape\color{gray},
        columns=fullflexible, keepspaces=true, breaklines=true,
        frame=single, framerule=0.3pt, rulecolor=\color{gray},
        captionpos=b, abovecaptionskip=4pt, xleftmargin=2pt, xrightmargin=2pt}
\renewcommand{\lstlistingname}{Código}
\renewcommand{\lstlistlistingname}{Índice de códigos}


\begin{document}
\crefname{table}{tabla}{tablas}
\Crefname{table}{Tabla}{Tablas}
\crefname{figure}{figura}{figuras}
\Crefname{figure}{Figura}{Figuras}
\crefname{code}{código}{códigos}
\Crefname{code}{Código}{Códigos}
\crefname{lstlisting}{código}{códigos}
\Crefname{lstlisting}{Código}{Códigos}

% -----------------------------------------------------------------------------
\chapter{Introducción\label{CAP:INTRODUCTION}}
  \section{Motivación\label{SEC:INTRO-MOTIVATION}}%
          {chapters/01-introduction/11-motivation}
  \section{Objetivos\label{SEC:INTRO-OBJECTIVES}}%
          {chapters/01-introduction/12-objectives}
  \section[Contexto y coordinación con el TFG de backend]%
          {Contexto del proyecto y coordinación con el \acs{tfg} de
           \textit{backend}\label{SEC:INTRO-CONTEXT}}%
          {chapters/01-introduction/13-context}
  \section{Estructura del documento\label{SEC:INTRO-STRUCTURE}}%
          {chapters/01-introduction/14-structure}

% -----------------------------------------------------------------------------
\chapter{Estado del arte\label{CAP:STATE-OF-THE-ART}}
  \section[Interfaces de plataformas de corrección]%
          {Interfaces de plataformas de corrección digital\label{SEC:SOTA-UI}}%
          {chapters/02-state-of-the-art/21-correction-uis}
  \section{Tecnologías de \textit{frontend}\label{SEC:SOTA-TECH}}%
          {chapters/02-state-of-the-art/22-frontend-tech}
  \section[Contratos de API y desarrollo dirigido por contrato]%
          {Contratos de \acs{api} y desarrollo dirigido por
           contrato\label{SEC:SOTA-CONTRACT}}%
          {chapters/02-state-of-the-art/23-api-contracts}
  \section[Agentes de programación basados en LLM]%
          {Agentes de programación basados en modelos de
           lenguaje\label{SEC:SOTA-AGENTS}}%
          {chapters/02-state-of-the-art/24-coding-agents}
  \section{Sistemas de diseño y accesibilidad\label{SEC:SOTA-DESIGN}}%
          {chapters/02-state-of-the-art/25-design-systems}

% -----------------------------------------------------------------------------
\chapter{Análisis\label{CAP:ANALYSIS}}
  \section[Análisis del backend existente]%
          {Análisis del \textit{backend} existente\label{SEC:AN-BACKEND}}%
          {chapters/03-analysis/31-backend-analysis}
  \section{Actores, roles y permisos\label{SEC:AN-ACTORS}}%
          {chapters/03-analysis/32-actors}
  \section{Requisitos funcionales del \textit{frontend}\label{SEC:AN-RF}}%
          {chapters/03-analysis/33-functional}
  \section{Requisitos no funcionales\label{SEC:AN-RNF}}%
          {chapters/03-analysis/34-nonfunctional}
  \section[Huecos del contrato]%
          {Huecos detectados en el contrato del
           \textit{backend}\label{SEC:AN-GAPS}}%
          {chapters/03-analysis/35-gaps}

% -----------------------------------------------------------------------------
\chapter{Diseño\label{CAP:DESIGN}}
  \section{Arquitectura general\label{SEC:DES-ARCH}}%
          {chapters/04-design/41-architecture}
  \section[Capa de datos, autenticación y simulación]%
          {Capa de datos, autenticación y simulación del
           \textit{backend}\label{SEC:DES-DATA}}%
          {chapters/04-design/42-data-layer}
  \section{Diseño de la interacción\label{SEC:DES-INTERACTION}}%
          {chapters/04-design/43-interaction}
  \section{Sistema de diseño visual\label{SEC:DES-VISUAL}}%
          {chapters/04-design/44-visual-system}
  \section{Anotación sobre el examen escaneado\label{SEC:DES-ANNOTATION}}%
          {chapters/04-design/45-annotation}

% -----------------------------------------------------------------------------
\chapter{Implementación\label{CAP:IMPLEMENTATION}}
  \section[Metodología: desarrollo asistido por agentes]%
          {Metodología: desarrollo asistido por agentes de
           programación\label{SEC:IMP-METHOD}}%
          {chapters/05-implementation/51-methodology}
  \section{Herramientas\label{SEC:IMP-TOOLS}}%
          {chapters/05-implementation/52-tooling}
  \section{Partes críticas\label{SEC:IMP-CRITICAL}}%
          {chapters/05-implementation/53-critical}
  \section[Coordinación y auditoría del trabajo paralelo]%
          {Coordinación y auditoría del trabajo en
           paralelo\label{SEC:IMP-PARALLEL}}%
          {chapters/05-implementation/54-parallel}
  \section[Uso de herramientas de IA]%
          {Uso de herramientas de inteligencia
           artificial\label{SEC:IMP-AI}}%
          {chapters/05-implementation/55-ai-use}

% -----------------------------------------------------------------------------
\chapter{Pruebas y validación\label{CAP:TESTING}}
  \section{Estrategia de pruebas\label{SEC:TEST-STRATEGY}}%
          {chapters/06-testing/61-strategy}
  \section{Pruebas de flujo\label{SEC:TEST-FLOWS}}%
          {chapters/06-testing/62-flow-tests}
  \section{Verificación determinista\label{SEC:TEST-AUDITS}}%
          {chapters/06-testing/63-audits}
  \section{Cobertura del contrato\label{SEC:TEST-COVERAGE}}%
          {chapters/06-testing/64-api-coverage}
  \section[Validación con el backend real]%
          {Validación con el \textit{backend} real\label{SEC:TEST-REAL}}%
          {chapters/06-testing/65-real-backend}

% -----------------------------------------------------------------------------
\chapter[Hacia la generación agéntica de frontends]%
        {Hacia la generación agéntica de \textit{frontends} a partir
         de un \textit{backend}\label{CAP:GENERALIZATION}}
  \section{El proceso como \textit{pipeline}\label{SEC:GEN-PIPELINE}}%
          {chapters/07-generalization/71-pipeline}
  \section[Qué es determinista y qué requiere razonamiento]%
          {Qué es determinista, qué requiere razonamiento y qué
           requiere a una persona\label{SEC:GEN-CLASSIFICATION}}%
          {chapters/07-generalization/72-classification}
  \section{Propuesta de arquitectura\label{SEC:GEN-PROPOSAL}}%
          {chapters/07-generalization/73-proposal}

% -----------------------------------------------------------------------------
\chapter{Conclusiones y trabajo futuro\label{CAP:CONCLUSIONS}}
  \section{Conclusiones\label{SEC:CONC-CONCLUSIONS}}%
          {chapters/08-conclusions/81-conclusions}
  \section{Trabajo futuro\label{SEC:CONC-FUTURE}}%
          {chapters/08-conclusions/82-future-work}

% -----------------------------------------------------------------------------
\appendix

\chapter{Inventario y cobertura de la API\label{AP:API}}%
        {appendices/A-api/A-overview}
\chapter{Manual de instalación y ejecución\label{AP:INSTALL}}%
        {appendices/B-install/B-overview}
\chapter{Registro de iteraciones\label{AP:ITERATIONS}}%
        {appendices/C-iterations/C-overview}
\chapter[Declaración detallada de uso de IA]%
        {Declaración detallada del uso de herramientas de
         IA\label{AP:AI}}%
        {appendices/D-ai-use/D-overview}
\chapter{Guía de estilo visual\label{AP:STYLE}}%
        {appendices/E-style/E-overview}

\end{document}
